% ===================================================================
%  TABLEAU N° 9 — La montagne. La forêt.
%
%  Traduction, faite en 2026, en francais standard du Quebec, sur
%  l'ido de texto/io. Regles de la colonne : voir 10-tableau-01.tex.
%  Deux scenes, deux numerotations : les cles portent c1 et c2.
%
%  LES QUATRE CHOIX PROPRES A CE TABLEAU :
%    (41) bohémien ............ Tsigane
%    (30) pic-vert ............ pic-bois
%    (52) piqueurs ............ chasseurs à cheval
%         myriamètres ......... dizaines de kilomètres
%
%  « BOHEMIEN » EST LE SEUL MOT DE LA COLONNE QU'ON REMPLACE POUR CE
%  QU'IL FAIT, ET NON POUR CE QU'IL DIT. L'ido ecrit « cigano » et le
%  fac-simile « bohemien » : les deux nomment la meme personne, et le
%  livret de 1926 n'y met aucune malveillance. Mais le mot francais a
%  servi depuis, et il designe aujourd'hui moins un peuple qu'un
%  cliche. « Tsigane », avec sa majuscule d'ethnonyme, nomme des gens.
%  Le reste de l'alinea ne bouge pas : l'homme conduit un ours dresse,
%  et c'est ce que la gravure montre.
%
%  « PIC-BOIS » EST LE MOT DU QUEBEC pour l'oiseau que la France
%  appelle pic-vert. Rien d'autre ne change dans la phrase : il
%  cherche toujours ses insectes dans la mousse.
%
%  « MYRIAMETRE » EST MORT, ET IL FALLAIT LE DIRE AUTREMENT. L'ido
%  ecrit « quaropa kilometri » — quatre fois dix kilometres, la vieille
%  unite de dix kilometres que le systeme metrique a laissee tomber.
%  Aucun lecteur de 2026 ne la connait. « Plusieurs dizaines de
%  kilometres » dit la meme distance sans faire semblant.
% ===================================================================

%%K t09-apar-1 apar
\VUsaut{4.0mm}
\VUtitre{15.0pt}{60}{TABLEAU N\textsuperscript{o} 9}
\VUsaut{3.0mm}

%%K t09-tit-1 sub
\VUcentre{12.0pt}{0}{\textit{Première scène.}}
\VUsaut{3.0mm}

%%K t09-tit-2 sub
\VUpk{11.4pt}{40}{La montagne. --- La ville d'eaux.}
\VUsaut{3.0mm}

%%K t09-c1-01-1 p
1. --- Voilà huit jours que je suis à Pratz. Je ne peux pas dire toute
l'admiration que j'ai éprouvée quand je me suis trouvé devant la merveilleuse
\VUgras{chaîne} \textsuperscript{(1)} des Pyrénées, dont la \VUgras{crête}
\textsuperscript{(2)} se découpe sur le ciel bleu comme un feston géant.

%%K t09-c1-02-1 p
2. --- Je ne m'imaginais pas que les \VUgras{sommets} \textsuperscript{(3)} des
Pyrénées atteignaient une \VUgras{hauteur} aussi grande, et je suis resté
stupéfait devant les magnifiques \VUgras{glaciers} \textsuperscript{(5)} et les
\VUgras{pics} \textsuperscript{(4)} neigeux sur lesquels roulent de si terribles
\VUgras{avalanches}.

%%K t09-tit-3 sub
\VUsaut{3.0mm}
\VUpk{10.2pt}{40}{Paysage de montagne.}
\VUsaut{3.0mm}

%%K t09-c1-03-1 p
3. --- Dès le lendemain de mon arrivée, je suis allé au vieux \VUgras{volcan}
\textsuperscript{(6)} éteint qui se trouve près de Pratz. Sur les bords arides
et dénudés du \VUgras{cratère}, on voit encore très bien les traces du passage
de la \VUgras{lave} qui a coulé, il y a plusieurs siècles, sur le \VUgras{flanc
de la montagne} \textsuperscript{(7)}. J'ai réussi à trouver, sur une des
\VUgras{pentes}, quelques fragments de cette lave.

%%K t09-c1-04-1 p
4. --- Pratz est situé sur la frontière, et la \VUgras{forteresse}
\textsuperscript{(8)} qui défend le \VUgras{col} \textsuperscript{(27)} séparant
la France de l'Espagne rappelle tout à fait ces vieux châteaux des bords du Rhin
qui semblent guetter une \VUgras{proie} du haut de leur rocher.

%%K t09-c1-05-1 p
5. --- Un énorme \VUgras{aigle} \textsuperscript{(9)}, qui a son nid non loin du
\VUgras{fort} \textsuperscript{(8)}, attire souvent mon attention. Cet
\VUgras{oiseau de proie}, dont le \VUgras{bec} est puissant, apporte souvent à
ses petits un agneau ou un chevreau qu'il tient dans ses \VUgras{griffes}. Ses
\VUgras{ailes} sont si grandes qu'il peut planer longtemps avec son lourd
fardeau.

%%K t09-tit-4 sub
\VUsaut{3.0mm}
\VUpk{10.2pt}{40}{L'excursionniste.}
\VUsaut{3.0mm}

%%K t09-c1-06-1 p
6. --- La plus agréable promenade, ici, est l'excursion à la \VUgras{cascade}
\textsuperscript{(10)} de Cau. La \VUgras{chute d'eau} tombe en tourbillonnant,
puis disparaît en un joli \VUgras{ruisseau} dans un \VUgras{bois de sapins}
\textsuperscript{(11)} situé à l'ouest de la ville. Nous avons monté à travers
ce bois jusqu'à l'\VUgras{observatoire météorologique} \textsuperscript{(12)},
et du haut du \VUgras{belvédère} nous avons découvert tout le \VUgras{panorama}
de la région. Le \VUgras{paysage} est merveilleux, et la \VUgras{nature} semble
prodiguer à ce pays tous les trésors dont elle dispose.

%%K t09-c1-07-1 p
7. --- Pour descendre, nous avons pris le \VUgras{funiculaire}
\textsuperscript{(13)} et nous nous sommes arrêtés à une petite
\VUgras{station}, près des \VUgras{ruines} \textsuperscript{(14)} d'un vieux
\VUgras{château féodal} autrefois habité par les seigneurs de Pratz.

%%K t09-c1-08-1 p
8. --- Pauvre château! Que doit-il penser en voyant à ses pieds la coquette
\VUgras{ville d'eaux} \textsuperscript{(15)}, devenue une \VUgras{station
thermale} à la mode?

%%K t09-c1-08-2 p
Le \VUgras{climat} y est si agréable, l'\VUgras{eau minérale} si efficace, que
la \VUgras{cure} qu'on vient faire au \VUgras{sanatorium} \textsuperscript{(17)}
est un vrai plaisir.

%%K t09-tit-5 sub
\VUsaut{3.0mm}
\VUpk{10.2pt}{40}{Une agréable ville thermale.}
\VUsaut{3.0mm}

%%K t09-c1-09-1 p
9. --- On a d'ailleurs tout fait pour rendre le séjour agréable. Un grand
\VUgras{viaduc} \textsuperscript{(16)} a été construit pour permettre d'établir
une voie ferrée; le \VUgras{parc} \textsuperscript{(18)} de
l'\VUgras{établissement thermal}, où l'on donne des \VUgras{concerts}, est
aménagé de façon charmante. De jolies \VUgras{villas} \textsuperscript{(19)} ont
été bâties autour du casino \textsuperscript{(21)}, et une \VUgras{chaussée}
\textsuperscript{(22)} très bien entretenue facilite grandement les
communications.

%%K t09-c1-10-1 p
10. --- Un simple \VUgras{talus} \textsuperscript{(23)} sépare la
\VUgras{route} \textsuperscript{(20)} de la \VUgras{vallée}
\textsuperscript{(25)} proprement dite, au fond de laquelle s'étale un joli
petit \VUgras{lac} \textsuperscript{(26)} qui alimente un \VUgras{ruisseau}
\textsuperscript{(24)} limpide.

%%K t09-c1-11-1 p
11. --- De l'autre côté de la vallée, on exploite une \VUgras{mine de fer}
\textsuperscript{(28)}. Je suis allé plusieurs fois voir les \VUgras{mineurs}
descendre dans le \VUgras{puits}, et je suis même entré dans la
\VUgras{galerie} où ils passent leur journée à extraire le \VUgras{minerai} qui
contient le \VUgras{métal}.

%%K t09-c1-12-1 p
12. --- Une importante \VUgras{fonderie} a été construite près de la mine pour
utiliser tout de suite les richesses qu'on en retire. Le \VUgras{haut fourneau}
\textsuperscript{(29)}, chauffé au \VUgras{charbon de terre} qui abonde ici,
permet d'obtenir rapidement et à bon marché de la \VUgras{fonte}.

%%K t09-c1-13-1 p
13. --- Non loin de la mine, on creuse une \VUgras{carrière}
\textsuperscript{(30)}. Ce n'est ni du \VUgras{granit}, ni du \VUgras{porphyre},
ni même du \VUgras{grès} que le vieux \VUgras{carrier} détache avec son
\VUgras{pic}, mais simplement de la \VUgras{pierre de taille} pour la
construction des maisons.

%%K t09-c1-14-1 p
14. --- Un des plus beaux ornements de ce pays est certainement le
\VUgras{sapin} \textsuperscript{(31)}. Partout, au fond d'un \VUgras{ravin}
\textsuperscript{(32)}, au bord d'un \VUgras{torrent} \textsuperscript{(33)},
d'un \VUgras{gouffre} \textsuperscript{(34)} ou d'un précipice, leurs fines
aiguilles mettent une nuance verte.

%%K t09-tit-6 sub
\VUsaut{3.0mm}
\VUpk{10.2pt}{40}{La randonnée.}
\VUsaut{3.0mm}

%%K t09-c1-15-1 p
15. --- Je profite de mes vacances pour faire de \VUgras{longues promenades}.
Les \VUgras{chemins} \textsuperscript{(35)} qui serpentent sur la montagne
permettent de parcourir, sans faire attention à la distance, plusieurs
\VUgras{kilomètres} et souvent plusieurs dizaines de kilomètres.

%%K t09-c1-15-2 p
Des \VUgras{bornes} \textsuperscript{(36)} et des \VUgras{poteaux indicateurs}
\textsuperscript{(37)} jalonnent les chemins, où l'on rencontre souvent un jeune
\VUgras{montagnard} \textsuperscript{(38)} avec sa \VUgras{besace}
\textsuperscript{(40)} au côté, portant une \VUgras{marmotte}
\textsuperscript{(39)}, ou bien quelque \VUgras{Tsigane} \textsuperscript{(41)}
dont la \VUgras{toque de fourrure} \textsuperscript{(42)} contraste étrangement
avec le vieux \VUgras{capuchon} \textsuperscript{(43)} déchiré. Il conduit par
une \VUgras{chaîne} un \VUgras{ours} \textsuperscript{(44)} dressé.

%%K t09-tit-7 sub
\VUpk{10.2pt}{40}{L'ascension.}
\VUsaut{3.0mm}

%%K t09-c1-16-1 p
16. --- Presque tous les jours, je vais avec un \VUgras{guide}
\textsuperscript{(48)} faire une \VUgras{ascension}, et je suis très fier quand,
le \VUgras{sac} \textsuperscript{(47)} au dos et l'« \VUgras{alpenstock} » à la
main, en vrai \VUgras{touriste} \textsuperscript{(46)}, j'escalade les
\VUgras{rochers} \textsuperscript{(45)}.

%%K t09-c1-17-1 p
17. --- Du haut de la montagne, les habitants de la plaine ne paraissent pas
plus gros que des fourmis; un \VUgras{troupeau de moutons}
\textsuperscript{(50)} paissant dans un \VUgras{pâturage} \textsuperscript{(49)}
semble un jouet d'enfant. Un \VUgras{pin} \textsuperscript{(51)}, ou même un
\VUgras{chalet} \textsuperscript{(52)}, n'a l'air que d'une tache sur la
verdure.

%%K t09-c1-18-1 p
18. --- Pendant ces excursions, nous rencontrons souvent dans le
\VUgras{sentier} \textsuperscript{(53)} un \VUgras{chasseur de chamois}
\textsuperscript{(54)} qui porte sur l'épaule la bête qu'il vient de tuer.

%%K t09-c1-19-1 p
19. --- J'ai mis dans mon herbier une branche de \VUgras{fougère}
\textsuperscript{(55)} très curieuse et un joli brin rose de \VUgras{bruyère}
\textsuperscript{(57)} que j'ai cueillis à l'entrée d'une petite \VUgras{grotte}
\textsuperscript{(56)} lors de ma dernière promenade.

%%K t09-tit-8 sub
\VUsaut{3.0mm}
\VUcentre{12.0pt}{0}{\textit{Deuxième scène.}}
\VUsaut{3.0mm}

%%K t09-tit-9 sub
\VUpk{11.4pt}{40}{La forêt. --- La chasse.}
\VUsaut{3.0mm}

%%K t09-c2-01-1 p
1. --- Hier, j'ai passé une journée délicieuse dans la forêt. L'activité qui
règne parmi les animaux et les \VUgras{insectes} \textsuperscript{(5)} qui
l'habitent m'a beaucoup intéressé.

%%K t09-c2-01-2 p
L'\VUgras{araignée} \textsuperscript{(1)} qui tisse sa \VUgras{toile}
\textsuperscript{(2)} entre deux branches pour attraper la \VUgras{mouche}
\textsuperscript{(3)} qui passe, l'\VUgras{escargot} \textsuperscript{(4)} qui
porte péniblement sa coquille, le lucane qui inspecte sa prise, la
\VUgras{fourmi} laborieuse qui s'agite auprès de la \VUgras{fourmilière}
\textsuperscript{(6)}, tous ces petits êtres allaient, venaient, travaillaient
avec une ardeur extraordinaire.

%%K t09-c2-02-1 p
2. --- Un \VUgras{serpent} \textsuperscript{(7)}, que j'ai d'abord pris pour une
\VUgras{vipère} mais qui n'était qu'une simple \VUgras{couleuvre}, s'était tapi
sous un tronc d'arbre et sortait de temps en temps sa petite tête fine, qui
semblait toujours prête à mordre. Devant moi, une grosse \VUgras{limace}
\textsuperscript{(8)} rampait lourdement après avoir dévoré une des délicieuses
petites \VUgras{fraises} \textsuperscript{(11)} qui poussent là en abondance.

%%K t09-c2-03-1 p
3. --- Le sol était tapissé de \VUgras{fleurs des bois} : des
\VUgras{primevères} \textsuperscript{(9)}, des \VUgras{pervenches}
\textsuperscript{(10)} et beaucoup d'autres plantes que je ne connaissais pas
fleurissaient au milieu des \VUgras{touffes d'herbe} \textsuperscript{(12)}.

%%K t09-c2-04-1 p
4. --- Dans cette région, la \VUgras{truffe} est presque aussi commune que le
\VUgras{champignon} \textsuperscript{(14)}, et un \VUgras{petit cochon}
\textsuperscript{(13)} appartenant à un garde forestier cherchait
gourmandement s'il n'en trouverait pas quelques-unes.

%%K t09-tit-10 sub
\VUsaut{3.0mm}
\VUpk{10.2pt}{40}{Le braconnage.}
\VUsaut{3.0mm}

%%K t09-c2-05-1 p
5. --- J'ai assisté à la \VUgras{fin} d'une scène qui m'a bien amusé. Grâce au
flair de son \VUgras{chien basset} \textsuperscript{(15)}, le \VUgras{garde
forestier} \textsuperscript{(16)} venait de surprendre un \VUgras{braconnier}
\textsuperscript{(22)} au moment où celui-ci se préparait à emporter le
\VUgras{gibier} \textsuperscript{(17)} qu'il avait tué. Aimant mieux abandonner
sa prise que de se laisser arrêter, le braconnier s'était enfui, et
\VUgras{lièvre} \textsuperscript{(18)}, \VUgras{biche} \textsuperscript{(19)},
\VUgras{chevreuil} \textsuperscript{(20)} et \VUgras{sanglier}
\textsuperscript{(21)} gisaient sur l'herbe, à la grande joie du garde
forestier.

%%K t09-c2-06-1 p
6. --- La variété des arbres que contient la forêt est étonnante. Les
\VUgras{hêtres} \textsuperscript{(23)} et les \VUgras{bouleaux}
\textsuperscript{(26)}, les \VUgras{pins}, qui produisent la \VUgras{résine}
\textsuperscript{(27)}, les \VUgras{chênes} \textsuperscript{(29)} et bien
d'autres encore mêlent leurs branchages.

%%K t09-c2-06-2 p
Le \VUgras{lierre} \textsuperscript{(24)}, qui s'accroche au tronc des grands
arbres, donne à la \VUgras{futaie} \textsuperscript{(25)} une verdure brillante
qui s'unit agréablement à la nuance plus claire et plus douce de la
\VUgras{mousse} \textsuperscript{(31)}, dans laquelle le \VUgras{pic-bois}
\textsuperscript{(30)} cherche des insectes.

%%K t09-tit-11 sub
\VUsaut{3.0mm}
\VUpk{10.2pt}{40}{Paysage de forêt.}
\VUsaut{3.0mm}

%%K t09-c2-07-1 p
7. --- Les vieux chênes m'ont charmé. Leurs grosses \VUgras{racines}
\textsuperscript{(32)} tordues, à demi sorties du sol, leur donnent un
splendide air de force et de vigueur, et je me demande comment le
\VUgras{propriétaire} \textsuperscript{(35)} peut se résoudre à les faire
abattre et à les livrer à la \VUgras{scie} \textsuperscript{(34)} du
\VUgras{bûcheron} \textsuperscript{(33)}. Les pauvres \VUgras{troncs}
\textsuperscript{(36)}, dépouillés de leur \VUgras{écorce}
\textsuperscript{(37)} ou fendus par la \VUgras{cognée} \textsuperscript{(38)}
du \VUgras{journalier} \textsuperscript{(39)}, me font de la peine.

%%K t09-c2-08-1 p
8. --- Tout près, un vieux \VUgras{charbonnier} \textsuperscript{(41)}
préparait avec sa \VUgras{pelle} une \VUgras{meule} \textsuperscript{(42)} de
charbon, et une vieille dame emportait un \VUgras{fagot} \textsuperscript{(40)}
de \VUgras{bois mort} fait des branchages des arbres abattus.

%%K t09-tit-12 sub
\VUsaut{3.0mm}
\VUpk{10.2pt}{40}{La chasse.}
\VUsaut{3.0mm}

%%K t09-c2-09-1 p
9. --- Je voulais aller me reposer quelques instants à la \VUgras{maison
forestière} \textsuperscript{(46)}, mais quand je suis arrivé auprès de
l'\VUgras{étang} \textsuperscript{(43)}, sur lequel nage un beau \VUgras{cygne}
\textsuperscript{(45)}, je me suis aperçu qu'une récente \VUgras{inondation}
\textsuperscript{(44)} avait changé le chemin en un véritable
\VUgras{marécage}. J'ai dû faire demi-tour et, en passant près du
\VUgras{cèdre} \textsuperscript{(49)}, des \VUgras{peupliers}
\textsuperscript{(48)} et de l'\VUgras{orme} \textsuperscript{(47)} qui sont en
bordure de la forêt, j'ai vu déboucher d'un \VUgras{taillis}
\textsuperscript{(54)} un magnifique \VUgras{cerf} \textsuperscript{(50)},
poursuivi par une \VUgras{meute} \textsuperscript{(51)} acharnée qu'excitait
encore le \VUgras{cor} \textsuperscript{(53)} des \VUgras{chasseurs à cheval}
\textsuperscript{(52)}. La pauvre bête était complètement épuisée. Elle est
venue tomber mourante à quelques pas de moi.

%%K t09-c2-10-1 p
10. --- Le fils du garde forestier, que le bruit de la \VUgras{chasse à courre}
avait attiré, est sorti de la maison et nous sommes retournés tous les deux dans
la forêt. Il avait vu une \VUgras{pie} \textsuperscript{(56)} sur une branche
d'un vieux chêne. Il pensait que son nid devait être caché dans le
\VUgras{feuillage} \textsuperscript{(58)} et il voulait le prendre.

%%K t09-c2-10-2 p
Leste comme un \VUgras{écureuil} \textsuperscript{(61)}, il a grimpé de
\VUgras{branche} \textsuperscript{(55)} en branche, mais il n'a rien trouvé et
n'a réussi qu'à faire envoler une vieille \VUgras{corneille}
\textsuperscript{(60)} perchée sur un \VUgras{rameau} \textsuperscript{(57)}.
\VUsaut{4.0mm}
\VUfilet{22mm}
